\section{Закон сохранения энергии}

\subsection{Работа силы}

\subsubsection{Работа}

\textbf{Работа} — скалярная физическая величина, характеризующая меру воздействия силы на тело, приводящего к его перемещению.

\textbf{Элементарная работа} силы \(\mathbf{F}\) на бесконечно малом перемещении \(d\mathbf{r}\):
\[
dA = \mathbf{F} \cdot d\mathbf{r} = F \cdot ds \cdot \cos\alpha
\]
где \(\alpha\) — угол между вектором силы и вектором перемещения, \(ds = |d\mathbf{r}|\).

\textbf{Полная работа} при перемещении материальной точки из положения 1 в положение 2 вдоль некоторой траектории \(L\) вычисляется как криволинейный интеграл II рода:
\[
A_{12} = \int\limits_{L_{1\to2}} \mathbf{F} \cdot d\mathbf{r}
\]

Если сила зависит только от координат \(\mathbf{F}(\mathbf{r})\) и траектория задана параметрически, то
\[
A_{12} = \int\limits_{t_1}^{t_2} \left[ F_x\frac{dx}{dt} + F_y\frac{dy}{dt} + F_z\frac{dz}{dt} \right] dt.
\]

В случае, когда сила \(\mathbf{F}\) постоянна по величине и направлению, а траектория прямолинейна,
\[
A = F s \cos\alpha.
\]

\subsubsection{Мощность}

\textbf{Мощность} — физическая величина, равная скорости совершения работы:
\[
P = \frac{dA}{dt}.
\]

Мгновенная мощность может быть выражена через силу и скорость:
\[
P = \mathbf{F} \cdot \mathbf{v} = F v \cos\alpha,
\]
где \(\mathbf{v}\) — скорость точки приложения силы, \(\alpha\) — угол между векторами силы и скорости.

Если мощность постоянна, то работа за конечный промежуток времени равна
\[
A = P \cdot \Delta t.
\]

\subsubsection{Теорема о кинетической энергии}

Рассмотрим материальную точку массой \(m\), движущуюся под действием силы \(\mathbf{F}\). По второму закону Ньютона:
\[
\mathbf{F} = m\mathbf{a} = m\frac{d\mathbf{v}}{dt}.
\]

Элементарная работа силы на перемещении \(d\mathbf{r} = \mathbf{v}dt\):
\[
dA = \mathbf{F} \cdot d\mathbf{r} = m\frac{d\mathbf{v}}{dt} \cdot \mathbf{v}dt = m\mathbf{v} \cdot d\mathbf{v}.
\]

Учитывая тождество \(\mathbf{v} \cdot d\mathbf{v} = \frac{1}{2}d(\mathbf{v}\cdot\mathbf{v}) = \frac{1}{2}d(v^2)\), получаем:
\[
dA = \frac{m}{2}d(v^2).
\]

Интегрируя от начального состояния 1 к конечному состоянию 2:
\[
A_{12} = \int\limits_{1}^{2} dA = \frac{m}{2}\int\limits_{v_1}^{v_2} d(v^2) = \frac{m}{2}(v_2^2 - v_1^2).
\]

Таким образом, работа всех сил, действующих на материальную точку, равна изменению её кинетической энергии:
\[\boxed{
A_{12} = \frac{mv_2^2}{2} - \frac{mv_1^2}{2} = \Delta E_k
}\]

Эта формула представляет собой \textit{теорему об изменении кинетической энергии}: работа равнодействующей всех сил, приложенных к телу, равна изменению его кинетической энергии.

\subsection{Потенциал}

\subsubsection{Консервативные силы}

Силы называются \textbf{консервативными} \textit{(или потенциальными)}, если их работа не зависит от формы траектории, а определяется только начальным и конечным положениями тела:
\begin{enumerate}
    \item Работа по любой замкнутой траектории \( C \) равна нулю:
    \[ \oint_C \mathbf{F} \cdot d\mathbf{r} = 0 \]
    \item Примерами являются: сила тяжести, сила упругости, кулоновская сила.
\end{enumerate}

\subsubsection{Потенциальная энергия}

\textbf{Потенциальная энергия} $U$ --- это скалярная физическая величина, равная энергии взаимодействия тел или частей тела между собой, определяемая их взаимным расположением в пространстве:

\begin{enumerate}
    \item Значение потенциальной энергии зависит от выбора \textit{нулевого уровня} (точки, где $U=0$).
    \item Работа консервативных сил равна убыли потенциальной энергии системы, то есть изменению потенциальной энергии со знаком «минус»:
    \[
    A_{12} = -\Delta U = -(U_2 - U_1) = U_1 - U_2
    \]
\end{enumerate}

Виды потенциальной энергии:
\begin{itemize}
    \item \textbf{В поле тяжести Земли}:
    \[
    U = mgh
    \]
    где $h$ --- высота над выбранным нулевым уровнем.
    
    \item \textbf{Упруго деформированного тела} (пружины):
    \[
    U = \frac{kx^2}{2}
    \]
    где $k$ --- жесткость, $x$ --- величина деформации.
    
    \item \textbf{Гравитационного взаимодействия}:
    \[
    U = -G\frac{M m}{r}
    \]
    
    \item \textbf{Электростатического взаимодействия} (зарядов $q_1, q_2$):
    \[
    U = k\frac{q_1 q_2}{r}
    \]
\end{itemize}

\subsubsection{Градиент}
В каждой точке пространства, где действует консервативная сила, можно определить значение потенциальной энергии $U(x, y, z)$. Такая функция называется \textit{потенциальным полем (потенциалом)}.

Сила связана с потенциальной энергией через пространственную производную — \textbf{градиент}:
\[
\mathbf{F} = -\text{grad}\, U = -\nabla U
\]

\textbf{Оператор градиента} ($\nabla$, «набла») в декартовых координатах имеет вид:
\[
\nabla = \mathbf{i}\frac{\partial}{\partial x} + \mathbf{j}\frac{\partial}{\partial y} + \mathbf{k}\frac{\partial}{\partial z}
\]

Таким образом,
\[
\mathbf{F} = -\nabla U = -\left( \frac{\partial U}{\partial x}\mathbf{i} + \frac{\partial U}{\partial y}\mathbf{j} + \frac{\partial U}{\partial z}\mathbf{k} \right).
\]

Вектор $\text{grad}\, U$ направлен в сторону наиболее быстрого \textit{возрастания} энергии, а сила $\mathbf{F}$ направлена в сторону её наиболее быстрого \textit{убывания} (поэтому в формуле стоит знак «минус»).

\textbf{Эквипотенциальные поверхности} — это поверхности, на которых потенциальная энергия постоянна:
\[
U(x, y, z) = \text{const}
\]

Свойства эквипотенциальных поверхностей:
\begin{enumerate}
    \item В каждой точке вектор градиента $\nabla U$ (а значит, и сила $\mathbf{F}$) направлен \textbf{перпендикулярно} к эквипотенциальной поверхности.
    
    \item Работа силы $\mathbf{F}$ при перемещении тела вдоль эквипотенциальной поверхности равна нулю, так как
    \[
    dA = \mathbf{F} \cdot d\mathbf{r} = -\nabla U \cdot d\mathbf{r} = 0,
    \]
    поскольку $d\mathbf{r}$ лежит в касательной плоскости к поверхности $U=\text{const}$, а $\nabla U$ ортогонален этой поверхности.
    
    \item Эквипотенциальные поверхности \textbf{не пересекаются}. Через каждую точку поля проходит только одна эквипотенциальная поверхность.
\end{enumerate}

\subsection{Закон сохранения энергии}
Полная механическая энергия системы $E$ — это сумма её кинетической энергии $K$ и потенциальной энергии $U$:
\[
E = K + U
\]

По \textbf{теореме об изменении механической энергии}, изменение полной механической энергии системы равно работе внешних сил:
\[
\Delta E = A_{\text{внешн.}}
\]

По \textbf{закону сохранения энергии}, в изолированной системе, где отсутствуют внешние силы, полная механическая энергия остается постоянной:
\[
E = K + U = \text{const} \implies K_1 + U_1 = K_2 + U_2
\]